% Section 1.1, from lectures/L01/appendix/a_compilation.md.

\section{Compiling and Executing Code in Linux}
\label{c1:sec:compilation}\label{c1:app:a}
If you are using a Linux environment, you can use \code{make} to compile and execute your code.
This section describes how to install \textbf{WSL} in order to build and run code in a Linux
environment on a Windows computer.

\subsection{Installation}
\label{c1:sec:wsl}
WSL (\emph{Windows Subsystem for Linux}) allows running a Linux distribution in a terminal
environment directly inside Windows, without using a virtual machine. This book uses the Linux
distribution \textbf{Ubuntu}.

WSL is used here to:
\begin{itemize}
  \item Build, store, and run code examples.
  \item Compile and execute C and C++ programs.
  \item Work in a Linux environment using standard development tools such as \code{gcc},
    \code{g++}, \code{make}, and \code{cmake}.
\end{itemize}

\subsubsection{Installing WSL and Ubuntu}
\begin{enumerate}
  \item Open \textbf{Windows PowerShell as Administrator}.
  \item Run the following commands:
\end{enumerate}

\begin{shell}
wsl --install
wsl --set-default-version 2
\end{shell}

\noindent Alternatively, Ubuntu can be installed directly using:

\begin{shell}
wsl --install -d Ubuntu
\end{shell}

\note If the system asks for a restart, restart the computer before continuing.

Verify the installation with:

\begin{shell}
wsl --status
\end{shell}

\noindent Ensure that \textbf{Default Version: 2} is displayed.

\subsubsection{If Ubuntu was not installed automatically}
\begin{enumerate}
  \item Open \textbf{Microsoft Store}.
  \item Search for \code{Ubuntu}.
  \item Select the latest LTS version (currently \code{Ubuntu 24.04 LTS}).
  \item Click \textbf{Install}.
\end{enumerate}

\subsubsection{Starting Ubuntu}
\begin{enumerate}
  \item Type \code{Ubuntu} in the Windows search bar.
  \item Launch the distribution.
\end{enumerate}
At the first startup:
\begin{itemize}
  \item Choose a short and clear username (for example your first name).
  \item Choose a password.
  \item Store the password safely.
\end{itemize}
\note The password is used only inside the Linux environment.

\subsubsection{Installing basic development tools}
After Ubuntu has started, update the system and install the necessary tools:

\begin{shell}
sudo apt -y update
sudo apt -y upgrade
sudo apt -y install build-essential cmake
\end{shell}

\noindent This installs:
\begin{itemize}
  \item \code{gcc} and \code{g++}.
  \item \code{make}.
  \item Various development libraries.
  \item \code{cmake}.
\end{itemize}
Verify the installation:

\begin{shell}
gcc --version
g++ --version
\end{shell}

\subsection{Creating a new program}
\label{c1:sec:newprogram}
For each program, create a new directory, for example \file{example-dir}:

\begin{shell}
mkdir example-dir
cd example-dir
\end{shell}

\noindent Inside this directory, create a file called \file{Makefile}:

\begin{shell}
touch Makefile
\end{shell}

\subsection{A simple Makefile}
\label{c1:sec:makefile}
In this Makefile, add the following content:

\begin{makecode}
# Build and run the application as default.
default: build run

# Build the application.
build:
	@g++ main.cpp -o main -Wall -Werror -std=c++17

# Run the application.
run:
	@./main

# Clean the application.
clean:
	@rm -f main
\end{makecode}

\note
\begin{itemize}
  \item The indentation under the targets \code{build}, \code{run}, and \code{clean} must consist of
    \textbf{tabs}, not spaces.
  \item This Makefile builds an executable named \code{main}.
  \item Regarding the compiler flags:
  \begin{itemize}
    \item \code{-Wall} enables most compiler warnings.
    \item \code{-Werror} converts warnings into errors, which helps prevent subtle bugs.
    \item \code{-std=c++17} sets the C++ version to C++17.
  \end{itemize}
  \item To see the compilation commands executed by \code{make}, remove the \code{@} prefix.
\end{itemize}

\subsection{A simple Makefile with parameters}
\label{c1:sec:makeparams}
It is also possible to use parameters to make the Makefile easier to maintain, especially when the
number of source files increases:

\begin{makecode}
# Target application.
TARGET := main

# C++ compiler.
CXX_COMPILER := g++

# C++ compiler flags.
CXX_FLAGS := -Wall -Werror -std=c++17

# Source files.
SRC_FILES := main.cpp

# Build and run the application as default.
default: build run

# Build the application.
build:
	@$(CXX_COMPILER) $(SRC_FILES) -o $(TARGET) $(CXX_FLAGS)

# Run the application.
run:
	@./$(TARGET)

# Clean the application.
clean:
	@rm -f $(TARGET)
\end{makecode}

\subsection{Tips}
\label{c1:sec:tips}
For compilation to work correctly:
\begin{itemize}
  \item Place all header and source files (\code{.hpp} and \code{.cpp} files) directly in this
    directory.
  \item Make sure to include all source files (\code{.cpp} files) between \code{g++} and \code{-o}
    in the build target.
  \item For simple examples, all source files can be placed in the same directory.
\end{itemize}

\subsubsection{Embedded project structure}
In larger embedded projects, source files are typically organized in separate directories such
as:

\begin{console}
include/
    driver/
        gpio.hpp
source/
    driver/
        gpio.cpp
    main.cpp
Makefile
\end{console}

\noindent The Makefile can then be extended to support these directories:
\begin{itemize}
  \item The \file{include} directory is included via the compiler flags.
  \item The source files are listed one per line using line continuation (\code{\\}).
\end{itemize}

\begin{makecode}
# Target application.
TARGET := main

# C++ compiler.
CXX_COMPILER := g++

# C++ compiler flags.
CXX_FLAGS := -Wall -Werror -std=c++17 -Iinclude

# Source files.
SRC_FILES := source/driver/gpio.cpp \
                source/main.cpp \

# Build and run the application as default.
default: build run

# Build the application.
build:
	@$(CXX_COMPILER) $(SRC_FILES) -o $(TARGET) $(CXX_FLAGS)

# Run the application.
run:
	@./$(TARGET)

# Clean the application.
clean:
	@rm -f $(TARGET)
\end{makecode}

\subsection{Compiling a test program}
\label{c1:sec:testprogram}
You can compile and run your program using the following command:

\begin{shell}
make
\end{shell}

\noindent You can also build the program without running it:

\begin{shell}
make build
\end{shell}

\noindent You can run the program without compiling first:

\begin{shell}
make run
\end{shell}

\noindent You can remove compiled files:

\begin{shell}
make clean
\end{shell}

\noindent Create and run the following test program in a file named \file{main.cpp}:

\begin{cppcode}
#include <cstdio>

int main()
{
    std::printf("Hello to C++!\n");
    return 0;
}
\end{cppcode}

\noindent Run the program using:

\begin{shell}
make
\end{shell}

\noindent Expected output:

\begin{console}
Hello to C++!
\end{console}

\subsection{Notes}
\label{c1:sec:compilenotes}
\begin{itemize}
  \item We begin by including the standard header \code{<cstdio>}, which is the C++ version of the
    C header \code{<stdio.h>} and contains functionality for input and output operations.
  \item Most functionality in the C++ standard library exists inside the namespace \code{std}
    (\emph{standard}).
  \item A namespace works somewhat like a directory in a filesystem and helps prevent naming
    conflicts.
  \item When using something from the standard library we therefore prefix it with \code{std::},
    for example:
  \begin{itemize}
    \item \code{std::uint8_t}
    \item \code{std::size_t}
    \item \code{std::vector}
  \end{itemize}
  \item \code{std::printf} is the C++ version of the C function \code{printf} and is used to print
    text to the terminal.
\end{itemize}
